\tzpanchor{0138}
\tzpentryheader{0138}{S\textsuperscript{3} Universe Topology}

\begingroup\small
\begingroup\scriptsize
\setlength{\tabcolsep}{4pt}
\renewcommand{\arraystretch}{0.92}
\noindent\begin{tabularx}{\linewidth}{@{}p{0.24\linewidth}p{0.36\linewidth}X@{}}
\textbf{UUID} & \multicolumn{2}{l@{}}{\tzpcode{fdd04fc0-6355-5f03-a623-f537dc58a64c}}\\
\textbf{PiID} & \tzpcode{7253594081284} & \textbf{TZPi}\quad \tzpcode{4}\quad \textbf{PiPE}\quad \tzpcode{145}\\
\textbf{RTEID} & \textbf{Poloidal $\theta$}\quad \tzpcode{535.8846 rad} & \textbf{Toroidal $\phi$}\quad \tzpcode{00.8830 rad}\\
\textbf{TZPID} & \tzpcode{ID0138} & \textbf{Node Dependencies}\quad \tzpidref{0137}\\
\textbf{Registry} & \tzpcode{registry\_id\_0138} & \textbf{Scale}\quad simply connected finite universe model\\
\textbf{LOC Classification} & QA612 algebraic topology & QB981 cosmology\\
\textbf{arXiv Classification} & Primary: gr-qc \quad Cross-list: astro-ph.CO, math-ph & \\
\textbf{Primary Support} & Hopf Fibration Statement & Finite Volume Law\\
\textbf{Closure Support} & Photon Return Time & \\
\end{tabularx}\par
\endgroup
\endgroup

\tzpsection{Canonical Statement}
Rather than being infinite and flat, the cosmos wraps around itself like a three-sphere.

\tzpsection{Canonical Equation}
\[
\eta:S^3\to S^2,
\qquad
\mathrm{fiber}(\eta)=S^1
\]
\[
\mathrm{Vol}(S^3)=2\pi^2 R^3
\]
\[
t_{\mathrm{return}}=\frac{2\pi R}{c}
\]

\begin{minipage}[t]{0.49\linewidth}
\tzpsection{Canonical Symbol Key}
\begingroup
\small
\raggedright
\textbf{Source equation symbols.}\par
\begin{itemize}[leftmargin=1.35em,itemsep=1pt,topsep=1pt,parsep=0pt]
    \item $S^3$: three-sphere source domain.
    \item $S^2$: two-sphere target domain.
    \item $S$: source equation symbol.
    \item $\pi$: source equation symbol.
    \item $R$: global radius, boundary radius, or topological scale.
    \item $t_{\mathrm{return}}$: source equation symbol.
    \item $c$: speed of light or relativistic conversion factor.
\end{itemize}
\endgroup
\end{minipage}\hfill
\begin{minipage}[t]{0.47\linewidth}
\tzpsection{Wolfram Form}
\begingroup
\scriptsize
\raggedright
\textbf{Source Wolfram model.}\par
\begin{Verbatim}[fontsize=\tiny,breaklines=true,breakanywhere=true]
(* TZPID ID0138 generated source Wolfram model *)
ClearAll[K0138, Cr0138, T, R, A, W, I, N];
eq0138$1 = HoldForm[eta[S3] -> S2fiberetaS1];
eq0138$2 = HoldForm[Vol[S^3] == 2*Pi^2*R^3];
eq0138$3 = HoldForm[tReturn == (2*Pi*R)/c];

K0138 = HoldForm[{eq0138$1, eq0138$2, eq0138$3}];

N = "normalized TRAWIN closure object for S\textsuperscript{3} Universe Topology";
I = "Photon Return Time integration/comparison layer";
W = "Finite Volume Law wave or operator carrier";
A = "Hopf Fibration Statement admissible action/amplitude condition";
R = "Finite Volume Law rotation/curl preservation layer";
T = "Hopf Fibration Statement local source condition";

Cr0138[expr_] := HoldForm[N[I[W[A[R[T[expr]]]]]]];

Result0138 = Cr0138[K0138];
\end{Verbatim}
\endgroup
\end{minipage}

\par\vspace{0.45\baselineskip}
\noindent\begin{minipage}[t]{\linewidth}
\begingroup
\small
\raggedright
\textbf{TRAWIN Operator Key.}\par
\noindent\textbf{Time/localization operator:} $\mathsf{T}\!\left[\mathrm{local}\right]$ Hopf Fibration Statement local source condition.\par
\noindent\textbf{Rotation/curl operator:} $\mathsf{R}\!\left[\nabla\times\right]$ Finite Volume Law rotation/curl preservation layer.\par
\noindent\textbf{Action/amplitude operator:} $\mathsf{A}\!\left[\mathrm{admissible}\right]$ Hopf Fibration Statement admissible action/amplitude condition.\par
\noindent\textbf{Wave/d'Alembertian operator:} $\mathsf{W}\!\left[\Box\right]$ Finite Volume Law wave or operator carrier.\par
\noindent\textbf{Integration operator:} $\mathsf{I}\!\left[\int\right]$ Photon Return Time integration/comparison layer.\par
\noindent\textbf{Normalization operator:} $\mathsf{N}\!\left[\mathrm{normalized}\right]$ normalized TRAWIN closure object for S\textbackslash{}textsuperscript\{3\} Universe Topology.\par
\endgroup
\end{minipage}

\clearpage
\noindent\textbf{[ID 0138] S\textsuperscript{3} Universe Topology}\par
\vspace{0.35em}
\tzpsection{Derivation Layer}
\paragraph{Primary Equation.}
\[
\eta:S^3\to S^2,
\qquad
\mathrm{fiber}(\eta)=S^1
\]
\[
\mathrm{Vol}(S^3)=2\pi^2 R^3
\]
\[
t_{\mathrm{return}}=\frac{2\pi R}{c}
\]

\paragraph{Primary Derivation.}
The first line specifies a map, while the subsequent line provides its explicit form \( \qquad \). Verifying that the formula satisfies the declared mapping properties confirms that the map is well defined.
\paragraph{Equation Type.}
This statement is classified as a map relation.

\paragraph{Interpretation.}
This relation applies to the quantities defined above. In the TRAWIN reading, the equation functions as the mapping carrier for the entry’s information content. According to CHL, the derivation provides a certificate connecting the canonical proposition to its proof certificate.

\par\vspace{0.35em}
\noindent\textbf{Derivable Consequences.}\quad
\begingroup
\footnotesize
\raggedright
The canonical equation anchors S\textbackslash{}textsuperscript\{3\} Universe Topology by connecting the source condition to the derivation layer and proof certificate.
\par\endgroup

\tzpsection{Equation Support}
\noindent\textbf{1. Hopf Fibration Statement}\par
\[
\eta:S^3\to S^2,
\qquad
\mathrm{fiber}(\eta)=S^1
\]
\noindent This gives the fibre structure of the global topology.\par
\noindent\textbf{2. Finite Volume Law}\par
\[
\mathrm{Vol}(S^3)=2\pi^2 R^3
\]
\noindent This states the compactness of the universe model quantitatively.\par
\noindent\textbf{3. Photon Return Time}\par
\[
t_{\mathrm{return}}=\frac{2\pi R}{c}
\]
\noindent This closes the progression by tying the global radius to a direct observational-style timescale.\par

\clearpage
\noindent\textbf{[ID 0138] S\textsuperscript{3} Universe Topology}\par
\vspace{0.35em}
\tzpsection{Dictionary}
\begingroup\small
\tzpsection{Definition}
S  Universe Topology is a Axiom. The global topology of the universe is a three sphere S . The Hopf fibration  : S  to S  with fibre S  ensures finite volume and finite photon return time; the unive

% BEGIN TZP CONTEXTUAL MEANING
\tzpsection{Contextual Meaning}
\begingroup\footnotesize\setlength{\emergencystretch}{3em}\sloppy
\textbf{Formal statement:} The global topology of the universe is a three sphere S . The Hopf fibration  : S  to S  with fibre S  ensures finite volume and finite photon return time; the unive
\textbf{Contextual meaning:} Rather than being infinite and flat, the cosmos wraps around itself like a three sphere. Photons eventually return, and there are no holes or disconnected regions.
\textbf{Derivable consequences:} Source evidence records the entry as a reusable object whose consequences should be read from the matched formal statement and local proof route.
\textbf{Lemma support:} The global topology of the universe is a three sphere S . The Hopf fibration  : S  to S  with fibre S  ensures finite volume and finite photon return time; the unive\par
\endgroup
% END TZP CONTEXTUAL MEANING

\tzpsection{Technical Interpretation}
Compare observed cosmic web filament spacing and topology with the predicted values; analyze large scale structure surveys for flux tube features. ID 138: S  Universe Topology

\tzpsection{Equation Grounding}
The global topology of the universe is a three sphere S . The Hopf fibration  : S  to S  with fibre S  ensures finite volume and finite photon return time; the unive

\tzpsection{Interpretive Role}
Rather than being infinite and flat, the cosmos wraps around itself like a three sphere. Photons eventually return, and there are no holes or disconnected regions.
\endgroup

\tzpsection{TZP Thesaurus}
\begin{enumerate}[leftmargin=1.6em,itemsep=6pt,topsep=4pt]
\item \textbf{Hopf Fibration Statement}: The statement giving the fibre structure of the global topology.
\item \textbf{Finite Volume Law}: The quantitative statement of the compactness of the universe model.
\item \textbf{Photon Return Time}: The tying of the global radius to a direct observational-style timescale.
\end{enumerate}

\tzpsection{Encyclopedia Mathematica}
\begingroup\footnotesize\sloppy
The visual plate below is reserved for the source-specific equation and progression graphic for [ID 0138]. It is included only as a companion to the canonical equation, not as a substitute for the formal text.
\par\endgroup
\ifTZPDarkEdition
\tzpdictionaryvisualband{D:/TZPID/kdp_workflow/build/tikz_figures/ID0138_dictionary_figure_dark.pdf}
\fi

\clearpage
\begingroup\tzpsupportsetup
\noindent\textbf{\fontsize{10.8pt}{12.8pt}\selectfont [ID 0138] Constructive Logic and Proof Certificate}\par
\noindent\textbf{\fontsize{10pt}{12pt}\selectfont S\textsuperscript{3} Universe Topology}\par
\vspace{0.45em}
\begingroup
\footnotesize
\setlength{\parskip}{0.22em}
\setlength{\parindent}{0pt}
\noindent\textbf{Curry--Howard--Lambek Interpretation}\par
Under the Curry--Howard--Lambek correspondence, this registry entry is read as a typed proof
object: the stated source condition supplies the proposition, the derivation supplies the
morphism, and the proof certificate records the machine handles needed to check the obligation
in the formal registry.

\vspace{0.45em}
\noindent\textbf{CHL Proof}\par
\textbf{Statement:} rather than being infinite and flat, the cosmos wraps around itself like a three-sphere

\textbf{Logic Validation:}\\
\textbf{Proposition:} ``S³ Universe Topology is valid when the stated geometric or topological relation
preserves its invariant.''\\
\textbf{Proof:} The source statement names a geometric/topological carrier. The proof obligation is to
show that deformation, projection, or passage through the stated structure leaves the
relevant invariant or admissible region well defined.

\textbf{VERDICT:} \ensuremath{\checkmark}\ \textbf{TRUE} --- topological-invariance validation

\textbf{Formal Status:} \ensuremath{\checkmark}\ THEORETICAL -- source-backed CHL proof obligation

\textbf{Physical Status:} THEORETICAL -- requires model-specific empirical interpretation
\par\vspace{0.45em}
\noindent{\tiny\textbf{Isabelle/HOL.} \tzpplaincode{physics\_validity\_from\_model, external\_validity\_package, registry\_number\_matches, equation\_count\_matches}}\par
\ifTZPDarkEdition
\tzpproofvisualband{D:/TZPID/kdp_workflow/build/tikz_figures/ID0138_proof_certificate_figure_dark.pdf}
\fi
\endgroup
\endgroup
